% Aula 10 --- uma divisão em que o Gini e a entropia discordam, e a forma da curva que explica isso.
Ana é bióloga e usa uma árvore de classificação para distinguir mudas de duas espécies, a
espécie $1$ e a espécie $0$, a partir de características medidas no viveiro. Num nó com $12$
mudas, $6$ de cada espécie, a árvore compara duas divisões:
\begin{description}[leftmargin=1.4cm,style=nextline]
  \item[A] pela altura, separa dois grupos de $6$ mudas, um com $2$ da espécie $1$ e $4$ da
        espécie $0$ e outro com $4$ da espécie $1$ e $2$ da espécie $0$;
  \item[B] por uma mancha rara na folha, isola uma única muda, da espécie $0$, das outras $11$,
        que têm $6$ da espécie $1$ e $5$ da espécie $0$.
\end{description}
Numa região com proporção $p$ de mudas da espécie $1$, o índice de Gini é $G(p)=2p(1-p)$, a
entropia é $D(p)=-p\log_2p-(1-p)\log_2(1-p)$, com $0\log_20=0$, e o erro de classificação é
$E(p)=\min(p,1-p)$. Uma divisão em regiões com $n_1$ e $n_2$ mudas e proporções $p_1$ e $p_2$
tem impureza ponderada $\big[n_1I(p_1)+n_2I(p_2)\big]/(n_1+n_2)$, para cada medida $I$, e a
árvore escolhe a divisão de menor impureza ponderada.
\begin{itens}
  \item \pts{2,0} Calcule a impureza ponderada das divisões A e B pelas três medidas e diga
        qual divisão cada medida escolhe. \emph{Dados:} $\log_23\approx1{,}5850$,
        $\log_25\approx2{,}3219$ e $\log_211\approx3{,}4594$.
  \item \pts{2,0} Seja $I$ uma medida de impureza qualquer para duas classes, com
        $I(p)=I(1-p)$, $I(0)=0$ e $I(6/11)>0$. Mostre que A tem impureza ponderada
        menor que B se e somente se $I(1/3)/I(6/11)<11/12$. Calcule essa razão para as
        três medidas e explique que propriedade da curva da entropia a leva a escolher B.
\end{itens}

\begin{solucao}
O nó tem $12$ mudas, $6$ de cada espécie, e proporção $1/2$ da espécie $1$. Na divisão A, os
filhos têm $6$ mudas cada, com proporções $2/6=1/3$ e $4/6=2/3$ da espécie $1$. Na divisão B,
um filho tem $1$ muda, da espécie $0$, com proporção $0$, e o outro tem $11$, com proporção
$6/11$. As três medidas são simétricas, $I(p)=I(1-p)$, e se anulam em $p=0$.

\smallskip
\textbf{(a)} \emph{Gini.} Como $G(1/3)=G(2/3)=2\cdot\tfrac13\cdot\tfrac23=\tfrac49$ e $G(0)=0$,
\begin{align*}
  \text{A: }&\ \tfrac12\,G\big(\tfrac13\big)+\tfrac12\,G\big(\tfrac23\big)=\tfrac49\approx0{,}444,\\
  \text{B: }&\ \tfrac1{12}\,G(0)+\tfrac{11}{12}\,G\big(\tfrac6{11}\big)
  =\tfrac{11}{12}\cdot2\cdot\tfrac6{11}\cdot\tfrac5{11}=\tfrac5{11}\approx0{,}455 .
\end{align*}
\emph{Entropia.} Como $\log_2\tfrac32=\log_23-1$, vale
$D(1/3)=D(2/3)=\tfrac13\log_23+\tfrac23\log_2\tfrac32=\log_23-\tfrac23$. Em B, com $D(0)=0$ e
$\log_26=1+\log_23$,
\begin{align*}
  \text{A: }&\ \log_23-\tfrac23\approx1{,}5850-0{,}6667\approx0{,}918,\\
  \text{B: }&\ \tfrac{11}{12}\,D\big(\tfrac6{11}\big)
  =\tfrac1{12}\Big(6\log_2\tfrac{11}6+5\log_2\tfrac{11}5\Big)\\
  &\ =\tfrac1{12}\big(11\log_211-6\log_26-5\log_25\big)\approx\tfrac{10{,}934}{12}\approx0{,}911 .
\end{align*}
\emph{Erro de classificação.} Em A, cada filho erra $2$ das suas $6$ mudas, e a impureza
ponderada é $4/12=1/3\approx0{,}333$. Em B, o filho de uma muda não erra, e o de $11$ prediz a
espécie $1$ e erra as $5$ mudas da espécie $0$: $5/12\approx0{,}417$.
\begin{center}\small
\renewcommand{\arraystretch}{1.25}
\begin{tabular}{@{}lccc@{}}
\toprule
medida & A & B & escolha \\
\midrule
Gini & $0{,}444$ & $0{,}455$ & A \\
entropia & $0{,}918$ & $0{,}911$ & B \\
erro de classificação & $0{,}333$ & $0{,}417$ & A \\
\bottomrule
\end{tabular}
\end{center}
O Gini e o erro escolhem A; a entropia escolhe B. Como o nó é o mesmo nas duas divisões,
comparar as impurezas ponderadas é o mesmo que comparar as reduções: com $G=1/2$ e $D=1$ no
nó, o Gini cai $0{,}056$ em A e $0{,}045$ em B, e a entropia, $0{,}082$ em A e $0{,}089$ em B.

\smallskip
\textbf{(b)} Pela simetria, a impureza ponderada de A é
$\tfrac12I(1/3)+\tfrac12I(2/3)=I(1/3)$, e, como $I(0)=0$, a de B é
$\tfrac1{12}I(0)+\tfrac{11}{12}I(6/11)=\tfrac{11}{12}I(6/11)$. Dividindo por $I(6/11)>0$, o que
preserva o sentido da desigualdade,
\[
  I\big(\tfrac13\big)<\tfrac{11}{12}\,I\big(\tfrac6{11}\big)
  \iff \frac{I(1/3)}{I(6/11)}<\frac{11}{12}\approx0{,}917 .
\]
As três razões são
\begin{align*}
  \frac{G(1/3)}{G(6/11)}&=\frac{4/9}{60/121}=\frac{121}{135}\approx0{,}896,\\
  \frac{D(1/3)}{D(6/11)}&\approx\frac{0{,}918}{0{,}994}\approx0{,}924,\\
  \frac{E(1/3)}{E(6/11)}&=\frac{1/3}{5/11}=\frac{11}{15}\approx0{,}733,
\end{align*}
em que $D(6/11)\approx0{,}994$ sai do item (a), porque lá se obteve
$\tfrac{11}{12}D(6/11)\approx0{,}911$. Só a razão da entropia passa de $11/12$, de acordo com o
item (a).

A condição compara o valor de $I$ em $1/3$ com o valor em $6/11\approx0{,}545$, perto de $1/2$,
onde as três medidas atingem o máximo: $G(6/11)$ é $120/121\approx0{,}992$ do máximo do Gini, e
$D(6/11)\approx0{,}994$, do da entropia. A pergunta é, em essência, se $I(1/3)$ fica abaixo de
$11/12\approx0{,}917$ do máximo de $I$. Para o Gini, $G(1/3)/G(1/2)=8/9\approx0{,}889$; para a
entropia, $D(1/3)/D(1/2)\approx0{,}918$. Ao passar de $p=1/2$ para $p=1/3$, o Gini perde $11\%$
do seu máximo, e a entropia, só $8\%$: perto do máximo, a curva da entropia é mais achatada que a
do Gini. Por isso a entropia trata uma região com um terço de uma espécie como quase tão impura
quanto uma equilibrada, dá pouco valor à purificação parcial que A produz nos dois filhos e, em
comparação, mais valor à purificação completa de uma única muda, que é o que B faz.
\end{solucao}
